% brief/01_question.tex: the task, why it is hard, the approach, the status (Figure 1: the pipeline, one column).
% Wording from sections/01_introduction.tex, shortened. This is the only place that states the audit status; it also
% defines "lane", used by Figure 1's legend (short/diagram_pipeline.tex) and by Sections 4 and 5.
\section{The question}\label{sec:question}
The task, set by the principal investigator (PI), is to predict which hand brings a smartphone to the nose from the
motion sensors that the browser delivers to a web page at about 60~Hz~\cite{w3c2025deviceorientation}.

The answer is a percentage, \pright{}: the probability that the right hand moved the phone. It must be calibrated
(of the touches given 80\,\%, about 80\,\% should be right-hand touches) and hold for persons the model has never
seen.

The two hands make mirror-image movements, so simple descriptors such as speed and duration are much the same for
both; persons, phones and grips add their own variation. In production there is also no start cue: the page must
find the approach in a continuous stream and decide when it has seen enough.

Figure~\ref{fig:pipeline} shows the pipeline. Every evaluation leaves one person out (LOSO), because splits by
trial overstate the accuracy for new persons~\cite{saeb2017usecase}.

This version is preliminary: the analysis ran in lanes, one per model family or analysis step, and the final-model
lane, which selects, calibrates and exports the model, is a first pass, not yet independently audited. The full
report gives every decision with its evidence (Appendices~\fullref{app:data} to~\fullref{app:pd01}).

\begin{figure}[tbp]
\centering
% short/diagram_pipeline.tex: the overview diagram, shared by short.tex and brief.tex (\input it inside a figure).
% TikZ, drawn at a fixed size (about 7.4 cm wide, explicit 7.5/8.5-pt fonts), so it reads the same at column width in the
% brief and beside its caption in the short version. Numbers only through \R keys. Dashed boxes: the cue-free steps,
% simulated by lane m4 with its reference classifier yaw8; the final model has not yet run behind the detector.
\begin{tikzpicture}[
    node distance=2.4mm,
    stage/.style={draw=black!55, rounded corners=1.5pt, fill=black!3, text width=6.8cm, inner sep=2.4pt,
      align=left, font=\fontsize{7.5}{9}\selectfont},
    cuefree/.style={stage, dashed, fill=white},
    flow/.style={-{Latex[length=1.5mm,width=1.3mm]}, black!65, semithick},
    head/.style={font=\fontsize{8.5}{10}\selectfont\bfseries}]
  \node[stage] (sensors) {{\fontsize{8.5}{10}\selectfont\bfseries 1\enspace Phone sensors}\\
    accelerometer and gyroscope, read by the browser at about 60~Hz};
  \node[cuefree, below=of sensors] (detect) {{\fontsize{8.5}{10}\selectfont\bfseries 2\enspace Find the approach
    (no start cue)}\\ cut the stream into candidate movements; take the most likely approach if
    $P(\mathrm{approach}) \geq \R{m4/det_threshold}$, else wait};
  \node[stage, below=of detect] (std) {{\fontsize{8.5}{10}\selectfont\bfseries 3\enspace Standardise}\\
    gravity frame and heading-following body frame; \R{p1/n_norm} samples, 12 channels};
  \node[stage, below=of std] (model) {{\fontsize{8.5}{10}\selectfont\bfseries 4\enspace Model}\\
    1D CNN, trained only on other persons};
  \node[stage, below=of model] (calib) {{\fontsize{8.5}{10}\selectfont\bfseries 5\enspace Calibrated
    \pright{}}\\ Platt calibrator fitted inside each training fold};
  \node[cuefree, below=of calib] (stop) {{\fontsize{8.5}{10}\selectfont\bfseries 6\enspace Several attempts,
    then stop}\\ combine the logits, $\sigma\bigl(\sum_i z_i/\tau_k\bigr)$; answer once the result leaves
    $(1 - p_{\mathrm{stop}},\, p_{\mathrm{stop}})$, $p_{\mathrm{stop}} = \R{m4/p_stop}$, or after
    \R{m4/k_max} attempts\\[13pt]\mbox{}};
  % The answer bar inside box 6 (schematic, not to scale): left of the band LEFT, right of it RIGHT.
  \begin{scope}[shift={(stop.south west)}, xshift=2.5mm, yshift=1.4mm, font=\fontsize{7.5}{9}\selectfont]
    \fill[handleft!30]  (0,0) rectangle (1.25,0.30);
    \fill[black!10]     (1.25,0) rectangle (5.05,0.30);
    \fill[handright!30] (5.05,0) rectangle (6.3,0.30);
    \draw[black!50] (0,0) rectangle (6.3,0.30);
    \node[text=handleft]  at (0.625,0.15) {\sffamily\bfseries LEFT};
    \node[text=black!65]  at (3.15,0.15) {next attempt};
    \node[text=handright] at (5.675,0.15) {\sffamily\bfseries RIGHT};
  \end{scope}
  \foreach \a/\b in {sensors/detect, detect/std, std/model, model/calib, calib/stop}
    \draw[flow] (\a) -- (\b);
  \draw[flow] (stop.east) -- ++(3.2mm,0) |- (detect.east); % back to step 2 while P stays inside the band
  \node[below=1.2mm of stop.south west, anchor=north west, inner sep=0pt, font=\fontsize{7}{8.5}\selectfont,
    text=black!70] {dashed: simulated with lane m4's classifier yaw8};
\end{tikzpicture}%

\caption{From the phone's sensors to the answer shown. The dashed steps, which production needs without a start
cue, were simulated in lane m4 with a simpler reference classifier, yaw8 (Section~\ref{sec:results}), not yet with
the final model (Section~\ref{sec:production}).}
\label{fig:pipeline}
\end{figure}
